\section{Why FIFO discovery gives minimum edge counts}\label{maze-p09}

The start has distance zero. Processing a cell at distance zero can discover only cells at distance one. They are appended behind the current work. More generally, while processing a layer at distance d, new entries have distance d plus one and follow all entries already waiting from the current layer.

This ordering prevents a later longer route from claiming a cell before a shorter route would reach it. If an undiscovered cell had a route shorter than the distance being assigned, the predecessor on that shorter route would belong to an earlier layer and would already have discovered it. That contradicts its unseen state. The argument depends on unit-cost edges and on the queue actually being FIFO.

Two predecessors in the same layer can both offer the same shortest distance. Neighbor order determines which discovery occurs first. The program records that first predecessor and retains it. Processing the goal is the stopping rule chosen for this assignment; a different stopping policy would need corresponding changes to the publication contract, even if it could still produce a valid shortest path.
